Define the ideal norm and show some of its more important properties

Definition 6.1. The (absolute) ideal norm of an ideal \((0) \neq a \subset \mathcal{O}_K\) is
\[
\mathcal{N}(a) := (\mathcal{O}_K : a) \in \mathbb{N};
\]

it is finite, as seen in theorem 2.23.
Lemma 6.2. Let \(0 \neq \alpha \in \mathcal{O}_K\). Then we have \(\mathcal{N}((\alpha)) = |\text{N}_{K/\mathbb{Q}}(\alpha)|\).

Proof. This follows directly because \( \alpha \mathcal{O}_K \) is effectively the multiplication by \( \alpha \) matrix \(M_\alpha\) (with some choice of basis), so \(\text{N}_{K/\mathbb{Q}}(\alpha) = \det(A)\) and \(|\det A| = (\mathcal{O}_K : (\alpha))\) by the proof of Th. 2.23.

Theorem 6.3. Let \(a, b \neq (0)\) be ideals in \(\mathcal{O}_K\). Then we have \(\mathcal{N}(ab) = \mathcal{N} a \cdot \mathcal{N} b\).
Proof. Decompose \(a = p_1^{v_1} \cdots p_r^{v_r}\) with pairwise disjoint prime ideals 
\(p_i \neq (0)\), \(v_i \in \mathbb{N}\). By the Chinese remainder theorem (3.10 + 3.11):
\(
\mathcal{O}_K / a \cong \mathcal{O}_K / p_1^{v_1} \oplus \cdots \oplus \mathcal{O}_K / p_r^{v_r},
\)
by taking the absolute norm we get
\(
\mathcal{N} a = \# (\mathcal{O}_K / a) = \# (\mathcal{O}_K / p_1^{v_1}) \cdots \# (\mathcal{O}_K / p_r^{v_r}) = \mathcal{N}(p_1^{v_1}) \cdots \mathcal{N}(p_r^{v_r}).
\)
Claim: \(p^{v} \supseteq (p^{v})^{v}\) for any prime ideal \(p \neq (0)\), \(v \in \mathbb{N}\).

Consider \(\mathcal{O}_k ⊃ \mathfrak{p} ⊃ \mathfrak{p}^2 ⊃ \dots ⊃ \mathfrak{p}^v\), these inclusions are strict because of the unique
decomposition into prime ideals.

\(\pi^i/\pi^{i+1}\) is a \(\mathcal{O}_K/\mathfrak{p}\)-vector space, since \(\mathfrak{p} \cdot \pi^i = \pi^{i+1}\), so multiplication with 
\(\mathcal{O}\) is well-defined modulo \(\mathfrak{p}\).
Let \(a \in \pi^i \setminus \pi^{i+1}, b := (a) + \pi^{i+1}\), so \(\pi^i \supset b \supset \pi^{i+1}\). 
By the uniqueness of prime decomposition this implies \(b = \pi^i = (a) + \pi^{i+1}\), i.e. \(\pi^i/\pi^{i+1}\) is generated by \(a + \pi^{i+1}\), 
hence 1-dimensional. Using the third isomorphism theorem

\(A \supset B \supset C \ R\text{-modules} \implies A/B \cong (A/C)/(B/C)\),

we get
\(\mathcal{N}(\mathfrak{p}^v) = \#(\mathcal{O}_K/\mathfrak{p}^v) = \#(\mathcal{O}_K/\mathfrak{p}) \cdot \#(\mathfrak{p}/\mathfrak{p}^2) \cdots \#(\mathfrak{p}^{v-1}/\mathfrak{p}^v) = (\mathcal{N}\mathfrak{p})^v\).
since each \( \#(\mathfrak{p}^{i}/\mathfrak{p}^{i+1}) = \#(\mathcal{O}_K/\mathfrak{p})\).
The claim from the theorem now follows easily via the prime ideal decomposition of a and b.

State and prove the Minkowski bound for elements in ideals.

\( \textbf{Lemma 6.4.} \) 
Let \((0) \neq \mathfrak{a} \subset \mathcal{O}_K\) be an ideal. Then there is an \(0 \neq \alpha \in \mathfrak{a}\) with

\(|N_{K/\mathbb{Q}}(\alpha)| \leq \left(\frac{2}{\pi}\right)^s \sqrt{|d_K|} \mathcal{R}(\mathfrak{a})\)

Proof.
Let \( \varepsilon > 0 \). For each \( \tau \in G = \text{Hom}_{\mathbb{Q}}(K, \mathbb{C}) \) choose \( c_\tau \in \mathbb{R}_{>0} \) with 
\( c_\sigma = c_{\overline{\sigma}} \) for pairs of complex emb.
\[
\prod_{\tau \in G} c_\tau = \left(\frac{2}{\pi}\right)^s \sqrt{|d_K|} \mathcal{R}(\mathfrak{a}) + \varepsilon.
\]

Then theorem 5.4 guarantees a \( 0 \neq \alpha \in \mathfrak{a}\) with \( |\tau(\alpha)| < c_\tau \), \( \forall \tau \in G \). 
Hence
\[
(|N_{K/\mathbb{Q}}(\alpha)| = | \prod_{\tau \in G} \tau(\alpha) | < \prod_{\tau \in G} c_\tau = M + \varepsilon.
\]
As this works for any \( \varepsilon > 0 \) and \( |N_{K/\mathbb{Q}}(\alpha)| \) is always a positive integer (as \( \alpha \in \mathcal{O}_k \) ), there exists an \( \alpha \) with \( |N_{K/\mathbb{Q}}(\alpha)| \leq M \).



Show that the class number of a number field is finite

Let \( K \) be a number field. Then the class number \( h_K = \#\text{Cl}_K = \#(J_K/P_K) \) is finite.

Proof. Claim 1: \(\forall M > 0\) there are only finitely many ideals with \(\mathcal{N} \mathfrak{a} \le M\).
By multiplicativity of \(\mathcal{R}\) and since \(\mathfrak{R}(\mathfrak{a}) \in \mathbb{Z}_{\Sigma^2}\) for \(\mathfrak{a} \neq \mathcal{O}_K\), 
it is enough to show that only finitely many prime ideals (\mathfrak{p}) with (\mathcal{R} \mathfrak{p} \leq M) exist.

For each prime number \( p \), there are only finitely many prime ideals \(\mathfrak{p} \mid p\) 
(by uniqueness of prime ideal factorization of \( (p) \)).

For these \(p \supseteq \mathfrak{p}\) and \(\mathfrak{p} \nmid 1\), hence 
\(\mathbb{Z} \supseteq p \cap \mathbb{Z} \supseteq p \cdot \mathbb{Z}) \implies (\mathfrak{p} \supseteq \mathbb{Z} = p \cdot \mathbb{Z}\)
as \(p \cdot \mathbb{Z}\) is maximal in \( \mathbb{Z} \).

Consider the ring homomorphism \(\mathbb{Z} \to \mathcal{O}_K/\mathfrak{p}, \ n \mapsto n + \mathfrak{p}\), the kernel is 
\(\mathfrak{p} \cap \mathbb{Z} = p\mathbb{Z}\).

We obtain an injective field homomorphism

\(Z/p\mathbb{Z} \hookrightarrow \mathcal{O}_K/\mathfrak{p} \Longrightarrow p = \#(Z/p\mathbb{Z}) \le \#(\mathcal{O}_K/\mathfrak{p}) = \mathcal{R}(\mathfrak{p})\).

This shows that at most the prime ideals (\mathfrak{p}) dividing the prime numbers (p \leq M) can have (\mathcal{R} \mathfrak{p} \leq M).

Every ideal class in \(\text{Cl}_K\) contains an integral ideal \(\mathfrak{a}\) of norm \(\mathcal{R} \mathfrak{a} \leq M := (\frac{2}{\pi})^s \sqrt{|d_K|}\).

From a given ideal class, choose an arbitrary element; a fractional ideal \(\mathfrak{a}\). Choose \(0 \neq c \in \mathcal{O}_K\) so that 
\(c \mathfrak{a}^{-1} \subset \mathcal{O}_K\) (i.e. an integral ideal, see remark (ii) for Def. (3.12)).

By Lemma (6.4), \(\exists \alpha \in c\mathfrak{a}^{-1} \text{ with } |N_{K/\mathbb{Q}}(\alpha)| \le M \cdot \mathcal{R}(c \cdot \mathfrak{a}^{-1})\). 
Hence,
\(\mathcal{R}(\alpha c \mathfrak{a}^{-1}) = \mathcal{R}((\mathfrak{a})) \cdot \mathcal{R}(c^{-1}\mathfrak{a}^{-1}) \leq M\) 
\( \mathcal{R}(c^{-1}\mathfrak{a}^{-1}) = \mathcal{R}(c\mathfrak{a}^{-1})^{-1}\)
Furthermore, \(\alpha c \mathfrak{a}^{-1}\) is an integral ideal because \(\alpha \in c \mathfrak{a}^{-1} = (c^{-1} \mathfrak{a}^{-1}\), by definition.


Combining claim 1 and claim 2 yields the desired result.
